% Generated draft: requires mathematical and source review.
\documentclass[11pt]{article}
\usepackage[margin=25mm]{geometry}
\usepackage{fontspec}
\setmainfont{Latin Modern Roman}
\usepackage{amsmath,amssymb,mathtools,unicode-math}
\setmathfont{Latin Modern Math}
\usepackage{enumitem,needspace,xurl,hyperref}
\setlength{\parindent}{0pt}
\setlength{\parskip}{5pt}
\setlength{\emergencystretch}{3em}
\newcommand{\N}{\mathbb{N}}
\newcommand{\Z}{\mathbb{Z}}
\newcommand{\Q}{\mathbb{Q}}
\newcommand{\R}{\mathbb{R}}
\newcommand{\C}{\mathbb{C}}
\newcommand{\F}{\mathbb{F}}
\newcommand{\A}{\mathbb{A}}
\newcommand{\PP}{\mathbb{P}}
\newcommand{\E}{\mathbb{E}}
\newcommand{\eps}{\varepsilon}
\newcommand{\dd}{\,\mathrm{d}}
\begin{document}
\section*{UnsolvedMath ID: 20001870}
\textbf{Source problem code:} catalogue entry
\subsection{Definitions and mathematical statement}
Let \(F:TS^2\to[0,\infty)\) be a smooth strongly convex Finsler metric away from the zero section: \(F(x,av)=aF(x,v)\) for \(a>0\), and the fiber Hessian of \(F^2/2\) is positive definite. Reversibility \(F(x,-v)=F(x,v)\) is not assumed. A prime closed geodesic is a periodic geodesic not obtained by iterating a shorter one. Count orbits modulo time translation in the unit-speed geodesic flow; opposite orientations, when both are geodesics, represent distinct flow orbits.

The source conjecture, explicitly about irreversible metrics, is that the number \(N(F)\) of distinct prime closed geodesics satisfies
\[
N(F)=2\quad\text{or}\quad N(F)=\infty.
\]
The supplied research aid reports the stronger conclusion for every possibly irreversible Finsler metric, without bumpiness or transversality assumptions. It attributes the result to a 2026 JAMS theorem of Cristofaro-Gardiner, Hryniewicz, Hutchings and Liu. This is a supplied solved-in-literature assessment, not independent verification of that theorem or its exact hypotheses.

For the contact interpretation let \(S_FS^2=\{(x,v):F(x,v)=1\}\), diffeomorphic to \(\R\mathrm P^3\), and equip it with its Hilbert contact form. Its Reeb orbits are the unit-tangent lifts of geodesics. The aid notes that \(H^2(\R\mathrm P^3;\Z)\cong\Z/2\Z\) is torsion and reports that the stated Reeb theorem applies directly.

The supplied covering reduction is as follows. For a connected regular finite covering of degree \(D\), a simple downstairs periodic orbit whose monodromy has order \(o\) has \(D/o\) distinct simple lifts, each of period \(o\) times its downstairs period. In the universal double cover \(S^3\to\R\mathrm P^3\), let \(N_0\) and \(N_1\) count downstairs simple orbits contractible and noncontractible in the unit tangent bundle. The reported count is
\[
N_{\mathrm{up}}=2N_0+N_1.
\]
According to the aid, in the exactly-two case both unit-tangent orbits are noncontractible: \(N_0=0\), \(N_1=2\); each has one doubled-period lift and the lifted flow again has exactly two simple orbits. Both prime orbits are reported irrationally elliptic, meaning their transverse return-map eigenvalues are \(e^{\pm2\pi i\theta}\) with \(\theta\notin\Q\); all iterates are therefore nondegenerate, so the metric is bumpy. Conversely, a prime geodesic whose unit-tangent orbit is contractible forces infinitely many prime closed geodesics under these reported results. Contractibility here is in \(S_FS^2\), not in the simply connected base sphere.

The original historical support was a result of Hofer--Wysocki--Zehnder on contact \(S^3\), projected to the Finsler sphere under bumpiness and transverse intersections of stable and unstable manifolds of hyperbolic closed geodesics. It also cites Harris--Paternain and notes that the geodesics in Katok's example are elliptic. These earlier qualifications must not be mistaken for hypotheses of the aid's reported unrestricted conclusion.
\subsection{Short English statement}
The source conjectures that an irreversible Finsler two-sphere has exactly two or infinitely many prime closed geodesics. The supplied 2026 research aid labels this solved for all Finsler metrics and reports additional exact-two rigidity and covering consequences. These dated literature and orbit-counting claims remain subject to source and mathematical review.
\subsection{Sources}
\begin{itemize}
\item Keith Burns and Vladimir S. Matveev, \emph{Open problems and questions about geodesics}, extended report of the August 2010 International Workshop on Geodesics; \emph{Ergodic Theory and Dynamical Systems} 41 (2021), 641--684, labelled [BM21] in the supplied verification. Relevant passage: printed pages 3--4, Conjecture 2.2.1. AIM source: \url{https://aimath.org/pastworkshops/geodesicsproblems.pdf}.
\item The original passage identifies problem 15 from [\'Alvarez2006], cites H. Hofer, K. Wysocki and E. Zehnder [Ho-Wy-Ze2003], and A. Harris and G. Paternain [Har-Pat2008], and mentions Katok's example. Full bibliographic entries for these labels are not supplied.
\item The aid attributes the unrestricted theorem to Cristofaro-Gardiner, Hryniewicz, Hutchings and Liu, JAMS, 2026, without title, volume, pages or identifier.
\item Local registry \texttt{data/unsolvedmath/problems.json}, permanent ID 20001870, code \texttt{catalogue entry}; canonical \texttt{aim-geometry-notes.json}, \texttt{notes[207]}, source item 15; updated 14 August 2026. Statement status reconstructed unverified; the clean statement contains a corrupt ``Every S2 Finsler metric'' phrase, while the original specifies irreversible.
\item Attempt 1 is labelled full solution, solved, rigor proved, novelty confidence low. The research aid is machine generated, not peer reviewed, and requires independent verification before citation.
\end{itemize}
\end{document}
